\section{서론}
\label{sec:introduction}

% P1: 맥락, 관측 수단, 과제
이동하는 다수의 에이전트가 남기는 대규모 시공간 궤적은 사건 및 환경 기록과 결합될 때 에이전트 사이의 상호작용과 환경 변화까지 함께 분석할 수 있게 한다~\cite{Zheng15}.
배틀로얄 게임은 활동 영역이 줄어드는 환경에서의 생존 경쟁을 이러한 데이터로 대규모 관측할 수 있게 한다.
PUBG: BATTLEGROUNDS(이하 PUBG)에서는 최대 100명의 플레이어가 1인, 2인 또는 4인 팀으로 경쟁하며, 안전구역은 정해진 단계에 따라 줄어들고 구역 밖의 플레이어는 지속적인 피해를 입는다.
공개 텔레메트리는 모든 플레이어의 위치와 교전, 아이템, 구역 사건을 담고 있어~\cite{PUBGAPI}, 같은 규칙에서 반복되는 경기로 줄어드는 영역 아래의 생존을 반복 관찰할 수 있다.
새 안전구역이 공개되는 단계 경계마다 각 팀은 이동 경로와 교전 여부를 정해야 하므로, \rev{분석의 핵심 대상은 그 경계에서} 생존한 팀이 이어지는 한 단계 안에 탈락할 조건부 확률, 즉 다음 단계 탈락 위험이다.
이 양은 코치와 분석가의 경기 복기, 중계에서 주목할 팀의 선택, 개발사의 구역 설정 조정~\cite{PUBGBlueZone2026}에 활용될 수 있다.

% P2: 난점과 기존 설계의 한계
그러나 이 위험은 줄어드는 영역의 기하와 주변 팀과의 상호작용에 함께 달려 있고 그 수준이 단계마다 달라지며, 예측 시점 이후의 정보가 입력에 섞이면 결과가 누설된다.
기존 PUBG 예측 연구는 이 문제를 서로 다른 단위와 대상으로 \rev{다루는데,} Sagong 등~\cite{Sagong25a,Sagong25b}은 단계별 활동 점수로 개별 플레이어의 생존을 분류하고, KRAFTON의 예측기~\cite{KRAFTON}는 팀 간 Transformer로 탈락 순서의 상대위험을 학습하여 우승확률로 변환한다.
그러나 상대위험 점수는 정해진 다음 단계 안의 탈락확률과 다른 추정 대상이며, 그 학습 목적만으로는 단계마다 공통적인 위험 수준이 식별되지 않는다.
또한 두 설계는 분석 단위, 데이터, 평가 지표가 모두 달라, 어떤 설계 선택이 성능을 좌우하는지 보고된 결과만으로는 알 수 없다.

% P3: 생태학적 관점
줄어드는 영역에서의 생존은 서식지 손실 아래에서 집단의 존속을 다루는 생태학의 문제와 맞닿아 있다~\cite{Fahrig03}.
본 연구는 안전구역의 축소를 서식지의 단계적 감소로, 팀의 탈락을 집단의 소멸로 대응시켜, 다음 단계 탈락 위험을 단계를 시간 단위로 하는 이산시간 hazard로 정식화한다(\ref{sec:related:survival}절).
각 팀의 상태는 구역 경계에 대한 노출, 이동 부담~\cite{Nathan08}, 주변 경쟁 밀도, 남은 팀원 수로 구성한 서식지 맥락으로 표현한다.

% P4: 연구 질문과 접근
이 관점은 세 가지 질문을 제기한다.
\begin{itemize}
    \item \textbf{RQ1:}
    \rev{서식지 축소 관점의 네 요인, 즉 구역 경계에 대한 노출, 이동 부담, 주변 경쟁 밀도, 남은 팀원 수는 다음 단계 탈락 위험과 어떤 방향과 크기로 관련되며, 그 관계는 단계와 코호트에 걸쳐 일관되는가?}

    \item \textbf{RQ2:}
    단계 조건부 hazard 목적과 서식지 맥락 \rev{표현을 결합한 설계는 같은 데이터와 평가 조건에서 기존 설계의 적용판과 어떤 예측 성능 차이를 보이는가?}

    \item \textbf{RQ3:}
    입력 표현, 학습 목적, 단계별 기저위험, 모델 구조를 \rev{하나씩 바꿀 때 예측 성능은 얼마나 달라지며, 그 차이는 코호트마다 어떻게 다른가?}
\end{itemize}

이 질문에 답하기 위해 맵, 시점, 팀 규모로 나눈 12개 코호트의 공개 경기 텔레메트리에서 \rev{예측 시점 이후의 정보가 섞이지 않는 팀--단계 표본을 대규모로 구성하는 데이터 파이프라인을 구축한다.
이어서 두 기존 설계를 같은 표본과 평가 조건에 적용하여 본 연구의 설계와 비교하고, 네 설계 요소를 하나씩 바꾸어 각 요소의 조건부 기여를 평가한다.}

% 연구 질문에 대응하는 기여
본 논문의 기여는 세 질문에 대한 답으로 정리된다.
\begin{itemize}
    \item \textbf{RQ1에 대한 답.}
    다음 단계 탈락 위험을 서식지 축소 관점에서 정식화한 해석 가능한 모델을 제시한다.
    \rev{이동 부담과 주변 경쟁 밀도는 위험을 높이고 남은 팀원 수는 낮추는 방향이 평가 가능한 모든 코호트에서 일관되며, 구역 안쪽 깊이의 보호 효과는 구역이 작아질수록 커지고 경쟁 밀도의 효과는 단계가 진행될수록 작아진다.}

    \item \textbf{RQ2에 대한 답.}
    단계 조건부 hazard 목적과 서식지 맥락을 결합한 설계는 모든 코호트에서 두 기존 설계의 적용판보다 높은 예측 성능을 보이며, 특히 단계마다 달라지는 탈락 위험의 수준을 정확하게 나타낸다.

    \item \textbf{RQ3에 대한 답.}
    상대위험 기반 설계와의 차이는 주로 단계별 기저위험과 학습 목적이, 활동 점수 기반 설계와의 차이는 주로 서식지 맥락이 설명하며, 모델 구조가 더하는 이득은 실질적으로 작다.
\end{itemize}
